\section{Conclusion}
\label{sec:conclusion}

We have demonstrated that a St\"ormer--Verlet symplectic lattice with $1/\ln^2(t)$ adiabatic cooling---motivated by the Discrete Symplectic Cosmology framework---reproduces several key statistical features of the CMB temperature anisotropies.
Near-Gaussian one-point statistics emerge robustly from the symplectic dynamics, with an ablation study confirming that dissipative alternatives fail.
Acoustic-like oscillation peaks arise naturally from wave propagation on the lattice, and a three-parameter fit aligns all peak positions with a mock $\Lambda$CDM reference.
The $1/\ln^2(t)$ cooling law generates a finite acoustic horizon through asymptotic freeze-out, providing a mechanism for producing the sound horizon without abrupt decoupling.
Full-sky CMB maps from 3D lattice simulations show visual and statistical similarity to the Planck SMICA data.

These results establish the DSC lattice as an exploratory computational framework for generating CMB-like fluctuations from symplectic dynamics with spectrally-motivated cooling.
The differentiable implementation enables gradient-based inverse reconstruction, though a held-out test reveals that 2D CMB data alone cannot uniquely constrain the 3D structure (line-of-sight degeneracy); incorporating volumetric 3D anchors partially breaks this degeneracy, pointing toward multi-probe joint inference as the path to physically meaningful reconstruction.
Extending the reconstructed initial conditions forward through a nonlinear structure-formation phase reveals a V-shape thermodynamic transition---the $1/\ln^2(t)$ cooling naturally smooths primordial chaos into a uniform CMB without an inflaton field, while subsequent gravitational instability rebuilds complexity into a cosmic web.

Future work should pursue:
(i) a full Planck likelihood analysis (MCMC) to quantitatively constrain the lattice parameters;
(ii) multi-probe inverse reconstruction combining CMB with galaxy survey data (DESI, Euclid) to break the projection degeneracy;
(iii) extension to polarization and tensor modes;
(iv) connections between the lattice acoustic horizon and the physical sound horizon relevant to the Hubble tension.
